% Options for packages loaded elsewhere
\PassOptionsToPackage{unicode}{hyperref}
\PassOptionsToPackage{hyphens}{url}
\documentclass[
]{article}
\usepackage{xcolor}
\usepackage{graphicx}
\usepackage{amsmath,amssymb}
\setcounter{secnumdepth}{-\maxdimen} % remove section numbering
\usepackage{iftex}
\ifPDFTeX
  \usepackage[T1]{fontenc}
  \usepackage[utf8]{inputenc}
  \usepackage{textcomp} % provide euro and other symbols
\else % if luatex or xetex
  \usepackage{unicode-math} % this also loads fontspec
  \defaultfontfeatures{Scale=MatchLowercase}
  \defaultfontfeatures[\rmfamily]{Ligatures=TeX,Scale=1}
\fi
\usepackage{lmodern}
\ifPDFTeX\else
  % xetex/luatex font selection
\fi
% Use upquote if available, for straight quotes in verbatim environments
\IfFileExists{upquote.sty}{\usepackage{upquote}}{}
\IfFileExists{microtype.sty}{% use microtype if available
  \usepackage[]{microtype}
  \UseMicrotypeSet[protrusion]{basicmath} % disable protrusion for tt fonts
}{}
\makeatletter
\@ifundefined{KOMAClassName}{% if non-KOMA class
  \IfFileExists{parskip.sty}{%
    \usepackage{parskip}
  }{% else
    \setlength{\parindent}{0pt}
    \setlength{\parskip}{6pt plus 2pt minus 1pt}}
}{% if KOMA class
  \KOMAoptions{parskip=half}}
\makeatother
\setlength{\emergencystretch}{3em} % prevent overfull lines
\providecommand{\tightlist}{%
  \setlength{\itemsep}{0pt}\setlength{\parskip}{0pt}}
\usepackage{bookmark}
\IfFileExists{xurl.sty}{\usepackage{xurl}}{} % add URL line breaks if available
\urlstyle{same}
\hypersetup{
  hidelinks,
  pdfcreator={LaTeX via pandoc}}

\title{An Auditable Reliability-Regret Profile\\
for Retrospective ADMET QSAR Model Selection:\\
A Fixed-Benchmark Study Across 22 Endpoints}
\author{Huajian Jiang\\
College of Computing and Data Science,\\
Nanyang Technological University, Singapore\\
Correspondence: jian0157@e.ntu.edu.sg\\
ORCID: 0009-0001-3992-2286}
\date{}

\begin{document}

\maketitle

\subsection{Abstract}\label{abstract}

ADMET model selection is often based on one validation metric, although
predictive ranking, probability calibration, selective risk, residual
coverage and chemical support can disagree. We define an auditable
reliability-regret profile (RRP) that compares a family selected by a
designated validation metric with a retrospective, non-deployable
candidate-pool reference on prespecified reliability dimensions. The audit
covers 22 TDC ADMET endpoints, five model families and five seeds (550 fits),
with cleaning sensitivity, support-conditioned diagnostics and descriptive
source-transport analyses. Mean regret was 0.0325 for expected calibration
error, 0.0092 for Brier score, 0.0179 for coverage error, 0.0793 for relative
excess area under the risk-coverage curve and 0.1195 for applicability
alignment. Validation and reliability winners agreed in 4/13, 8/13, 2/9,
2/22 and 4/22 comparisons, respectively. Isotonic calibration reduced
expected calibration error on 13/13 classification endpoints, but Brier
improvements were uncertain. The planned confirmatory external matrix was
stopped by an outcome-blind power gate; the focused CYP2D6 analysis is
retrospective and exploratory because test-role outcomes were examined before
protocol freeze. The study supports multi-dimensional QSAR reliability
auditing within this fixed benchmark; it does not establish independent
external or temporal validity, deployment benefit, clinical utility or
acceptance.

\textbf{Keywords:} ADMET prediction; model selection; uncertainty
quantification; calibration; applicability domain

\subsection{1. Introduction}\label{introduction}

Computational QSAR prediction of absorption, distribution, metabolism,
excretion, and toxicity (ADMET) properties is used to prioritize
compounds and support decisions during drug discovery. Public benchmark
collections make it possible to compare models across endpoints, but the
resulting score depends on the endpoint, data version, split, and
evaluation metric. The Therapeutics Data Commons (TDC) was developed to
organize therapeutic machine-learning datasets, tasks, data splits, and
evaluation resources; its ADMET Group provides a common setting for
multi-endpoint assessment \cite{huang2021tdc}.

Because these are structure--property relationships, QSAR reliability also
depends on chemical support, assay heterogeneity, and endpoint-specific
calibration. A high predictive score alone therefore does not establish safe
extrapolation to compounds outside the observed chemical domain.

Many benchmark comparisons rank candidate models by a task-specific
predictive metric. Such a metric is useful for comparing predictive
ranking or error magnitude, but it does not by itself describe every
property relevant to using predictions. For classification, ranking
performance and probability calibration answer different questions: a
model can order compounds well while its predicted probabilities remain
poorly aligned with observed frequencies \cite{guo2017calibration}. In compound triage,
performance under selective review and the relation between structural
support and prediction error may also matter. These considerations
motivate evaluating a validation-selected model across multiple
reliability dimensions rather than treating its validation score as a
complete account of test behaviour.

We use these observations to define an auditable reliability-regret
profile (RRP). For each endpoint, the RRP records the loss incurred by the
family selected on validation for each prespecified reliability dimension,
relative to the most favorable member of the locked candidate pool on the
same held-out test partition. The reference is deliberately descriptive:
it exposes the dimensions that a single validation score leaves unresolved
without turning a test-set oracle into a deployable selector. The profile
therefore combines a computationally simple estimator with explicit
identity, split, seed, and provenance checks that can be rerun when the
candidate pool or data snapshot changes.

Recent work has already established several neighboring reliability questions. Hosni et al.~used a structural applicability-domain method with Venn--Abers or conformal uncertainty estimates to distinguish molecules for which uncertainty estimates were more versus less reliable under temporal, cluster, and industrial splits \cite{hosni2026}. Parrondo-Pizarro et al.~benchmarked chemical-distance, ensemble-variance, and error-model uncertainty metrics for ADME properties under temporal and controlled distribution shifts using separate training, calibration, and test roles \cite{parrondo2026}. Koleiev et al.~audited the top-ranked models across all 22 TDC ADMET leaderboards for reproducibility, leakage, and overfitting \cite{koleiev2026}. Castellanos et al.~analysed the prospective temporal-blind OpenADMET--ExpansionRx challenge and reported applicability-domain and error-tracking results; that work is a preprint and is cited here only as a complementary challenge report \cite{castellanos2026}. These studies delimit the present contribution: we do not introduce a new uncertainty estimator, calibration method, applicability-domain method, error model, temporal split, or reproducibility audit. The present contribution is the narrower reliability-regret selection audit and its evidence contract, which compares a validation-selected family with a retrospective candidate-pool reference across prespecified reliability dimensions.

Earlier benchmarking also examined data scarcity, structural distribution
shift, representation choice, data cleaning, cross-source transfer, and
structural-frontier controls in ADMET evaluation \cite{zhao2026reliability,kamuntavicius2025benchmarking,zheng2026frontier}.

Here, we conduct a retrospective audit of model selection and held-out
reliability across 22 TDC ADMET endpoints. For each endpoint, we select
one model family using the designated validation metric and compare that
choice with the candidate family that has the most favorable value on
each test-set reliability dimension. This test-set reference is a
descriptive, candidate-pool counterfactual; it is not an implementable
selection rule. The study therefore asks a bounded question: within this
fixed benchmark and candidate pool, how closely does selection by the
designated validation metric align with reliability measured on the
corresponding fixed test set? The dataset roles are separated in the
training code, but this analysis is not an independent external
confirmation: the same benchmark test endpoints have been examined
during prior development, and the primary TDC analysis used no
independent calibration cohort. A separate descriptive
transport diagnostic is reported in Section 2.8 and does not change this
limitation.

\subsection{2. Materials and Methods}\label{materials-and-methods}

\subsubsection{2.1. Study design and benchmark
tasks}\label{study-design-and-benchmark-tasks}

This retrospective benchmark audit covered 22 ADMET tasks in five
categories: absorption, distribution, metabolism, excretion, and
toxicity. Thirteen tasks were binary classification and nine were
regression. Endpoint metadata defined each task's designated metric. The
locked matrix comprised five model families and five random seeds (1-5)
for each endpoint, corresponding to 550 endpoint-model-seed fits. This primary TDC analysis evaluates model-family selection and
test-set reliability; it is not a prospective or clinical evaluation. A
separate source-specific external transport diagnostic is described
below and kept outside the primary TDC matrix.

The formal runs described in this draft used a PyTDC 1.1.15 data
snapshot staged on the remote compute system and audited as a distinct
data version. The earlier \texttt{full\_003} snapshot was unavailable
during this run and was not merged with the newer snapshot. Accordingly,
the data version in this manuscript must be reported as the audited
PyTDC 1.1.15 snapshot, not as an exact reproduction of the unavailable
historical run. Raw TDC records are not redistributed in the
public project package; data access remains subject to the source
dataset terms.

\subsubsection{2.2. Data identity and split
protocol}\label{data-identity-and-split-protocol}

For each endpoint, the analysis loaded the cached
\texttt{train\_val.csv} and \texttt{test.csv} files. For each seed, only
the \texttt{train\_val} partition was randomly divided into training and
validation subsets using scikit-learn \texttt{train\_test\_split} with
\texttt{test\_size=0.125}, shuffling enabled, and \texttt{random\_state}
equal to the task seed. Classification labels were used for
stratification; regression labels were not. This produced a nominal 7:1
training-to-validation split within the cached non-test rows. The cached
test membership and row order were retained and were not repartitioned.
The same task-seed split identities were used across all five model
families.

The split name \texttt{locked\_cache} denotes this explicit portable
protocol. It does not assert that the validation subset reproduces a
version-specific PyTDC default split or that validation is
scaffold-disjoint. The manuscript also does not label the cached test
set as scaffold-disjoint until the source-level provenance of that exact
snapshot has been verified. The test set was excluded from model
fitting, checkpoint selection, and calibration fitting. It was used for
final model evaluation and for constructing the descriptive
candidate-pool reliability reference. Because the same test endpoints
had been inspected during prior project development, the analysis is
retrospective and not a blinded confirmatory evaluation.

Two legacy pyridine N-oxide SMILES in the \texttt{Solubility\_AqSolDB}
source file required parser-compatibility translations under the current
RDKit runtime. The original source strings, labels, row indices, split
assignments, and prediction identities were retained. Only the private
strings passed to model parsers were translated under a finite, audited
mapping; no additional row was dropped, duplicated, relabelled, or
reassigned under this policy.

\subsubsection{2.3. Candidate model families and
training}\label{candidate-model-families-and-training}

The candidate pool comprised a random forest (RF), XGBoost, LightGBM, a
graph isomorphism network (GIN), and Chemprop's directed message-passing
neural network (D-MPNN) \cite{rogers2010ecfp,xu2019gnn,yang2019molecular,breiman2001rf,chen2016xgboost,ke2017lightgbm}.
The historical custom node-message-passing implementation was excluded
from the formal candidate pool and was not treated as Chemprop D-MPNN.

RF, XGBoost, and LightGBM used binary Morgan fingerprints generated with
radius 2 and 2,048 bits \cite{rogers2010ecfp}. RF used 500 trees, maximum depth 20,
minimum split size 5, and minimum leaf size 2; classification runs used
balanced-sub-sample class weights. XGBoost and LightGBM used 1,000
boosting iterations, maximum depth 6, learning rate 0.01, and
configuration values of 0.8 for row and feature subsampling. In the
locked LightGBM implementation, \texttt{subsample\_freq=0} disables row
bagging despite the \texttt{subsample=0.8} field, while feature
subsampling remains configured. No broad hyperparameter search
was conducted for this matrix; model configurations were fixed in the
project code and configuration files.

GIN operated on molecular graphs with 117-dimensional atom features
encoding atomic number, degree, formal charge, and hybridization. The
model used five trainable-epsilon GIN layers, hidden dimension 256, sum
pooling, and dropout 0.2. It was trained for up to 100 epochs with Adam
(learning rate 0.001; weight decay 0.00001), unweighted binary
cross-entropy for classification, and unscaled mean squared error for
regression. Early stopping was applied after 20 validation checks
without improvement \cite{xu2019gnn}.

Chemprop D-MPNN runs used Chemprop 2.2.1 with directed bond messages,
message hidden dimension 300, depth 3, zero dropout, normalized
aggregation with normalization constant 100, and a one-layer
feed-forward network with hidden dimension 300. Training was limited to
100 epochs with patience 20, two warm-up epochs, and batch size 64. The
per-run configuration records the selected validation metric and runtime
environment \cite{yang2019molecular}. These Chemprop settings refer to the formal
Chemprop implementation; parameters from the legacy custom
message-passing runner are not used to describe this model.

\subsubsection{2.4. Checkpointing and validation-based family
selection}\label{checkpointing-and-validation-based-family-selection}

Each endpoint-model-seed fit used its seed-specific training and
validation rows. The traditional estimators were fit once using their
locked configurations. For GIN and Chemprop, the checkpoint within a run
was selected using the task's configured validation metric and the
corresponding direction of optimisation. For these neural models,
classification checkpoint selection used AUROC or AUPRC; regression
checkpoint selection used MAE or Spearman correlation, except as noted
for Chemprop. Chemprop 2.2.1 did not track Spearman for checkpoint
selection, so its checkpoints on Spearman-designated tasks were selected
using validation MAE; final validation and test reporting for those runs
used the designated Spearman metric. This implementation-specific
difference is retained as a limitation of the candidate comparison.

For family selection, the designated metric was recomputed from each
seed's validation predictions, oriented so that larger values
represented better performance, and averaged arithmetically over the
five seeds. The family with the highest mean validation score was
selected for that endpoint. Validation rows were not pooled by row index
across seeds because each seed produced a different validation
partition. Near-tie membership and seed-bootstrap selection stability
were recorded separately from the winning-family rule.

\subsubsection{2.5. Test predictions and reliability
measures}\label{test-predictions-and-reliability-measures}

After family selection, the five seed-specific test predictions for each
eligible family were aligned by source sample identity. The per-compound
ensemble prediction was the arithmetic mean across seeds; ensemble
uncertainty was the sample standard deviation of those predictions. The
test labels were not used to select the validation winner. Test-set
reliability measures were calculated for every candidate family before
the selected-family value was compared with the most favorable candidate
value for the corresponding endpoint and reliability dimension.

Classification reliability included the Brier score, log loss, and
expected calibration error (ECE). ECE used ten equal-width probability
bins over {[}0, 1{]}. Regression reliability included empirical coverage
error for a nominal 90\% absolute-residual interval: a
finite-sample-corrected upper quantile of validation absolute residuals
(alpha=0.1) was estimated and applied to test predictions. Because
validation data were also used for checkpointing and family selection,
this interval is reported as a validation-residual coverage diagnostic
and not as an independent calibration result or an unconditional
conformal-coverage guarantee.

Selective-prediction behaviour was summarized with risk-coverage curves
that ordered compounds from lower to higher ensemble uncertainty.
Per-sample loss was Brier score and log loss for classification and
absolute error for regression. The area under the risk-coverage curve
(AURC) was compared with a loss-ordered oracle curve; excess AURC was
divided by the mean sample loss to obtain relative excess AURC. Ties in
uncertainty were handled by their expected within-tie ordering,
independent of input row order. The loss-ordered oracle is used only to
define a descriptive reference.

Applicability support was measured as each test compound's maximum
Tanimoto similarity to the endpoint's training compounds using radius-2,
2,048-bit Morgan fingerprints. Applicability alignment was the Spearman
correlation between this maximum similarity and negative absolute
prediction error; higher values indicate that more-supported test
compounds tended to have lower errors in this benchmark.

\paragraph{Reliability-regret profile.}
Let $e$ index endpoints, $m$ index the locked candidate families, and
$d$ index reliability dimensions. Let $V_{e,m}$ be the mean validation
score over the five seeds after orienting the metric so that larger is
better, and let $L_{e,m,d}$ be the corresponding test loss after
reorienting higher-is-better reliability measures as losses. The
validation-selected family is
\[
\hat m_e = \arg\max_m V_{e,m},
\]
with the frozen family order resolving exact ties. The dimension-specific
regret is
\[
R_{e,d}=L_{e,\hat m_e,d}-\min_m L_{e,m,d},
\]
and the range-normalized profile is
\[
\widetilde R_{e,d}=\frac{R_{e,d}}
{\max_m L_{e,m,d}-\min_m L_{e,m,d}+\epsilon},
\]
where $\epsilon$ prevents division by zero. The output is the vector
$\mathbf R_e=(R_{e,d})_d$ together with the candidate availability,
sample-identity audit, seed-level values, and endpoint-level uncertainty.
This estimator is a diagnostic of selection completeness: it does not
estimate a population effect, and the minimizing family is never used as a
future-data selection rule. Support-conditioned and selective-risk
analyses are attached to the same endpoint-level profile rather than
pooled across compounds.

For each endpoint and reliability dimension, the reliability regret was
the non-negative difference between the validation-selected family's
reliability loss and the lowest reliability loss observed among the
eligible candidate families on the test set. Higher-is-better dimensions
were reoriented as losses before this comparison. The test-set oracle is
an after-the-fact candidate-pool reference and cannot be used to select
a model for future compounds. Raw and range-normalized regret, tie-aware
summaries, and regret relative to uniform random selection were retained
in the analysis outputs. Dimension-specific practical tolerances and the
combined evidence gate are code-defined analysis criteria, not
prospectively registered hypotheses.

\subsubsection{2.6. Statistical summaries}\label{statistical-summaries}

The endpoint was the aggregation unit for cross-task summaries. For the
support-conditional analyses, model-task contrasts were first averaged
within endpoint; endpoint means were then used for the reported means,
percentile 95\% bootstrap intervals, sign-flip diagnostics, and
endpoint-positive counts. The \texttt{n\_pairs} field in the source
table counts eligible model-task contrasts before this endpoint
aggregation. Percentile 95\% intervals for mean regret summaries were
obtained by 5,000 bootstrap resamples of endpoint-level values;
validation-family selection stability used 1,000 resamples of the five
seeds. Monte Carlo sign-flip p-values used 20,000 repetitions with Holm
adjustment across reliability dimensions. Since regret is non-negative
by construction relative to a test-set oracle, these p-values are
retained as descriptive diagnostics and are not used by the evidence
gate. TDC endpoints may share compounds or source lineage, so
endpoint-bootstrap intervals are descriptive summaries of this fixed
task set and should not be interpreted as population-level inference
over independent assays.

\subsubsection{2.7. Reproducibility and data
access}\label{reproducibility-and-data-access}

The formal analysis uses explicit split identifiers, seed-specific
prediction files, and task/model run manifests. Prediction audits check
sample identity, row order, and label alignment, while per-run records
preserve model identifiers, seeds, and artifact hashes. GIN and Chemprop
formal fits used NVIDIA A40 GPUs; the Chemprop manifests record Python
3.12.10, Chemprop 2.2.1, PyTorch 2.7.1+cu128, RDKit 2026.09.1, NumPy
2.5.4, and CUDA 12.8. The formal five-model ledger was executed in the
tracked remote research environment. The local Mac was used for document
preparation, figure rendering, an independent replay audit of the 330
traditional-model artifacts, and the descriptive external diagnostics,
including the 120-fit traditional lane and the B3DB inference lanes; the Mac
replay is a verification lane and is not substituted for the primary ledger.
The Mac replay
environment was macOS 26.6 arm64 with Python 3.11.15, scikit-learn
1.9.1, RDKit 2026.09.1, XGBoost 3.2.0, and LightGBM 4.7.0. The retained
local manifests do not contain a complete package-version inventory for
the formal remote traditional CPU lane, so those versions are not
inferred or claimed; the full inventory is recorded in
\texttt{docs/CPU\_ENVIRONMENT\_INVENTORY\_20261021.md}.
The source package does not contain raw TDC records or model
checkpoints, so released prediction views and hashes alone cannot
independently reproduce label-dependent metrics; exact replay requires
an authorized matching TDC snapshot and the documented split protocol.

\subsubsection{2.8. Descriptive external transport
diagnostic}\label{descriptive-external-transport-diagnostic}

The separate diagnostic covered eight cohort-source combinations aligned
to four TDC classification tasks. For each cohort, RF, XGBoost, and
LightGBM were trained only on the corresponding locked TDC development
rows (\texttt{train\ +\ valid}), with five task seeds. External records
were divided into a calibration role and a sealed test role using the
frozen identity assignment; Platt and isotonic transforms were fit on
the calibration role, and the sealed test role was used only for
evaluation. We report source-by-model metrics as five-seed means and
standard deviations, without pooling cohorts or treating seeds as
independent external replicates.

Locked-checkpoint inference added 20 GIN prediction files across four
source-endpoint cohorts and 10 Chemprop prediction files across two
hERG source cohorts. External labels were not used for fitting or
checkpoint selection in these neural-model diagnostics; they remain
descriptive transport observations.

An independent B3DB diagnostic was kept in a separate ledger. After exclusion
of identities overlapping the locked TDC snapshot and the frozen v3 external
cohorts, the candidate contained 1,078 connectivity-deduplicated identities,
with 390 records assigned to calibration and 688 to sealed test. The B3DB
ledger comprised RF, XGBoost, and LightGBM fits for five seeds (15 files) and
locked GIN and Chemprop inference for five seeds each (10 files). B3DB labels
were not used for model fitting or checkpoint selection.

The planned confirmatory external analysis required a separate 150-fit
matrix and a prespecified outcome-blind power gate. That gate did not pass,
and the B3DB candidate's outcome-blind log-loss power screen also failed
(power 0.118 for a +0.010 log-loss effect under the prespecified simulation),
so neither descriptive diagnostic is a substitute for the unrun confirmatory
analysis.

\subsubsection{2.9. Validation-only ExpansionRx replication}\label{validation-only-expansionrx-replication}

To assess whether validation-stage family selection was sensitive to a
source-specific replication, we ran a post-hoc analysis on eight regression
endpoints from the public ExpansionRx replication package. The analysis
included five candidate families (RF, XGBoost, LightGBM, GIN, and Chemprop
D-MPNN), five seeds per family, and therefore 200 model slots. Models were
trained and scored only on the replication package's development and
validation roles. The validation-reliability selector (VRRS) combined equal
within-endpoint rank scores for validation MAE and the Spearman association
between maximum Morgan similarity to the seed-specific training set and
negative absolute validation error. Exact ties were resolved alphabetically.

The selector protocol was written after the replication execution and is
therefore explicitly post-hoc rather than preregistered. The selection table
and its leave-one-seed stability analysis were written before any test-target
directory was opened. No ExpansionRx test targets were available to this
analysis, so this replication cannot estimate unseen-test performance or
validate the selector as a deployable rule.

\subsubsection{2.10. Exploratory CYP2D6 diagnostic}\label{exploratory-cyp2d6-diagnostic}

A focused CYP2D6 diagnostic was run for the endpoint
\texttt{CYP2D6\_Veith}, and the model family was XGBoost selected solely from
the locked TDC validation table, using five seeds and two source cohorts (EPA
Tox21 and NCATS/NCGC). Each model was fit only on the locked TDC
train-plus-validation rows. However, the same source-specific test-role
metrics had already been reported in a diagnostic before the focused protocol
was frozen, and all ten later rerun prediction files were byte-identical to
the earlier candidate files. The test roles therefore were not untouched for
confirmatory inference.

For audit traceability, the diagnostic estimand was uncalibrated test log loss
minus Platt-calibrated test log loss, summarized as the equal-weight mean of
the two source-level means across seeds. Input, hash, and arithmetic checks
pass, but they do not repair prior test-role exposure. The result is retained
as exploratory only and supports no confirmatory CYP2D6 calibration or
transport claim.

\subsubsection{2.11. AI-assisted preparation and
accountability}\label{ai-assisted-preparation-and-accountability}

Generative AI assisted with study-code development and review, literature and
evidence organisation, document preparation, and language editing. Before
submission, the author will identify each tool and version and specify its
actual role in accordance with the journal policy. AI systems were not treated
as authors or scientific authorities. The human author will independently
verify the code, analyses, references, and manuscript and retains responsibility
for the accuracy, integrity, and interpretation of the submitted work.

\subsection{3. Results}\label{results}

\subsubsection{3.1. Analysis coverage and artifact
completion}\label{analysis-coverage-and-artifact-completion}

The final matrix covered all 22 endpoints, all five candidate families,
and five seeds per family, corresponding to 550 fit slots. The
reliability-regret manifest reported exact task, model, seed, and
sample-identity coverage, with 79 applicable endpoint-dimension cells
and 31 task-type-inapplicable cells out of 110 grid cells. The final
figure suite included the five-model matrix for the
conditional-reliability and calibration analyses. A separate
paired cleaning sensitivity analysis contained 120 runs across eight
endpoints and three traditional model families; its source, run, and
prediction audits passed, and the summary gate was
\texttt{SENSITIVITY\_STABLE}.

\subsubsection{3.2. Reliability regret after validation-based
selection}\label{reliability-regret-after-validation-based-selection}

Validation-selected models showed non-zero reliability regret across
multiple dimensions. Mean ECE regret was 0.0325 (95\% CI, 0.0148-0.0539)
across 13 classification endpoints, with a practically meaningful regret
on 9/13 endpoints. Mean Brier regret was 0.0092 (95\% CI,
0.0034-0.0169), with meaningful regret on 6/13 endpoints. For the nine
regression endpoints with coverage diagnostics, mean coverage-error
regret was 0.0179 (95\% CI, 0.0099-0.0258), with meaningful regret on
6/9 endpoints. Across all 22 endpoints, mean relative excess-AURC regret
was 0.0793 (95\% CI, 0.0401-0.1251) and mean applicability-alignment
regret was 0.1195 (95\% CI, 0.0739-0.1770); both dimensions had
meaningful regret on 15/22 endpoints, and the tie-robust counts were
12/22 for each dimension.

\subsubsection{3.3. Agreement between validation selection and
reliability
winners}\label{agreement-between-validation-selection-and-reliability-winners}

The validation-selected family and the test-set reliability winner often
differed. Exact winner agreement was 4/13 endpoints for uncalibrated
ECE, 8/13 for uncalibrated Brier score, 2/9 for coverage error, 2/22 for
selective-prediction gain, and 4/22 for applicability alignment. The
corresponding mismatch fractions were 69.2\%, 38.5\%, 77.8\%, 90.9\%,
and 81.8\%, respectively. These values are candidate-pool comparisons
made after observing the fixed test set; they do not define a deployable
test-set selection rule.

\subsubsection{3.4. Reliability conditional on chemical
support}\label{reliability-conditional-on-chemical-support}

At the primary maximum-similarity threshold of 0.5, the
outside-to-inside log-error ratio was positive for 20/21 eligible
endpoints, with a mean of 0.2654 (95\% CI, 0.1960-0.3390). The source
table contains 105 eligible model-task contrasts for this measure;
estimates and intervals are based on endpoint means after averaging
available model-task contrasts within each endpoint. The
uncertainty-correlation advantage for supported compounds was positive
for 14/21 endpoints, with a mean of 0.0850 (95\% CI, 0.0420-0.1314). The
selective-gain contrast was heterogeneous and centered below zero (mean
-0.0407; 95\% CI, -0.0939 to 0.0005). For the 12 classification
endpoints with the calibration contrast, the Brier gap was positive for
11/12 endpoints (mean 0.0398; 95\% CI, 0.0241-0.0567), whereas the ECE
gap was near zero (mean -0.0009; 95\% CI, -0.0147 to 0.0118).

\subsubsection{3.5. Calibration and cleaning
sensitivity}\label{calibration-and-cleaning-sensitivity}

Isotonic calibration reduced ECE on all 13 classification endpoints,
with a mean ECE reduction of 0.0221 (95\% CI, 0.0155-0.0294; exploratory
Holm-adjusted sign-flip p=0.0010). The corresponding Brier reduction was
small and uncertain (mean 0.0010; 95\% CI, -0.0015 to 0.0037). Platt
scaling showed no similarly consistent endpoint-level ECE or Brier
improvement in these summaries. In the cleaning sensitivity
analysis, the selected traditional-model winner changed for BBB\_Martins
and PPBR\_AZ, but no task met the gate's material-effect criterion; the
overall status was \texttt{SENSITIVITY\_STABLE} and the prediction audit
status was \texttt{GO}.

\subsubsection{3.6. Descriptive external transport
diagnostic}\label{descriptive-external-transport-diagnostic-1}

The first transport diagnostic completed 120 traditional-model fits across
eight cohort-source combinations, three model families, and five seeds.
Locked-checkpoint inference added 20 GIN prediction files and 10 Chemprop
prediction files on selected source-endpoint cohorts. Models were fit only on
TDC development rows; external calibration and test roles were kept separate,
and calibration transforms were fit on the calibration role alone. A separate
B3DB diagnostic completed 15 traditional-model fits and 10 locked neural
prediction files on 1,078 connectivity-deduplicated identities, with 390
calibration and 688 sealed-test records. The supplementary source-by-model
tables report metrics as means and standard deviations across the five seeds,
without pooling cohorts or treating seeds as independent external replicates.
The outcome-blind power screen for the planned confirmatory external analysis
did not pass, and the B3DB candidate power screen was also below the
prespecified threshold; the 150-fit broad confirmatory matrix therefore
remained at \texttt{0/150 POWER\_FAIL\_HOLD}. The focused CYP2D6 diagnostic
is reported separately in Section~3.8 because its test-role outcomes were
examined before protocol freeze; it is excluded from confirmatory external
evidence.

\subsubsection{3.7. Validation-only ExpansionRx replication}\label{validation-only-expansionrx-results}

The ExpansionRx replication completed all 200 planned model slots across
eight regression endpoints, five model families, and five seeds. VRRS selected
Chemprop D-MPNN for \texttt{caco2\_efflux}, \texttt{caco2\_papp}, and
\texttt{logd}; RF for \texttt{hlm\_clint} and \texttt{mlm\_clint}; GIN for
\texttt{ksol}; and LightGBM for \texttt{mbpb} and \texttt{mppb}. The VRRS
choice differed from the validation-MAE-only choice on 6/8 endpoints. The
leave-one-seed stability diagnostic agreed with the full VRRS choice in 33/40
comparisons (82.5\%). Because the test-target directory was not provided, no
post-freeze external test comparison was performed. The result supports only
validation-stage selector stability within this source-specific replication.

\subsubsection{3.8. Exploratory CYP2D6 diagnostic}\label{exploratory-cyp2d6-results}

The CYP2D6 diagnostic completed 10/10 XGBoost seed packages across the EPA
Tox21 and NCATS/NCGC source cohorts. The exploratory equal-weight source-level
paired log-loss improvement was 0.1630 (95\% source-stratified bootstrap CI,
0.1521--0.1739; 12,545 test rows); the pooled row-weighted summary was 0.1585.
These metrics are reported only to preserve the diagnostic record: they were
examined before the focused protocol freeze, and the later prediction files
were byte-identical to the earlier candidate outputs. They are not
confirmatory external evidence.

\subsubsection{3.9. Figures and source
files}\label{figures-and-source-files}

Figure 1 summarizes validation-selected reliability regret across the
five reliability dimensions; the colour scale is normalized regret and
the markers encode the locked practical and tie-robust thresholds.
Figure 2 shows conditional reliability contrasts by maximum training-set
similarity at the primary threshold of 0.5. Figure 3 reports
endpoint-level calibration gains and failures for isotonic and Platt
scaling. Editable SVG, publication PDF, PNG, and source-data files are
retained in the corresponding \texttt{results/figures/} directories,
with hashes recorded in the figure manifests.

\begin{figure}[htbp]
\centering
\includegraphics[width=0.98\linewidth]{figures/fig01_validation_selected_reliability_regret.png}
\caption{Validation-selected reliability regret across the five audited
dimensions and 22 ADMET endpoints. Values and endpoint coverage are linked to
the final reliability-regret manifest.}
\label{fig:reliability-regret}
\end{figure}

\begin{figure}[htbp]
\centering
\includegraphics[width=0.98\linewidth]{figures/fig02_conditional_reliability.png}
\caption{Conditional reliability contrasts by maximum training-set similarity
at the primary threshold of 0.5. The contrasts are descriptive endpoint-level
summaries and are not prospective performance guarantees.}
\label{fig:conditional-reliability}
\end{figure}

\begin{figure}[htbp]
\centering
\includegraphics[width=0.98\linewidth]{figures/fig03_calibration_gains_failures.png}
\caption{Endpoint-level calibration gains and failures for isotonic and Platt
scaling. The calibration analysis is internal to the fixed benchmark and does
not establish an independent calibration cohort.}
\label{fig:calibration-gains}
\end{figure}

\subsection{4. Discussion}\label{discussion}

This audit shows that a validation metric can identify a reasonable
predictive family without identifying the family that is best on every
reliability dimension of the fixed held-out set. The largest and most
recurrent mismatches involved selective prediction and applicability
alignment, while Brier-score agreement was higher than agreement for
those dimensions. The result is therefore a warning against treating one
validation score as a complete proxy for how uncertainty, calibration,
and chemical support behave across an endpoint.

The conditional analysis gives a more specific interpretation.
Lower-support compounds generally had larger errors and weaker
support-error relationships, but the selective-prediction contrast was
not uniformly favorable. Calibration was also dimension-specific:
isotonic scaling improved ECE consistently in this benchmark, while
Brier improvements were not stable. These patterns argue for reporting
several reliability dimensions together and for distinguishing a
descriptive benchmark audit from an operationally validated
model-selection procedure.

The methodological contribution is the reliability-regret profile and its
reproducible audit contract: validation selection, dimension-specific test
losses, support-conditioned diagnostics, seed aggregation, and artifact
lineage are reported as one separable object. This is a computational
evaluation framework rather than a new molecular architecture, chemical
split, or claim of superiority over the closest recent ADMET reliability
studies. The full-text comparison in
\texttt{docs/NEAR\_NEIGHBOR\_COMPARISON\_20261020.md} and the 2026 full-text audit record the narrower estimand and the split/model differences that limit direct comparison. In particular, the present RRP is complementary to applicability-domain and uncertainty-under-shift studies, ADME uncertainty benchmarks, TDC reproducibility audits, and the prospective OpenADMET--ExpansionRx challenge; it does not supersede or reproduce those studies.

The validation-only ExpansionRx replication adds a source-specific stability
check for the selection procedure, but it does not change the estimand of the
primary audit. Its post-hoc VRRS rule and the absence of test targets mean that
the result is best treated as an implementation and sensitivity diagnostic,
not as a new model-selection algorithm or an external generalization result.

\subsection{5. Limitations}\label{limitations}

First, the main analysis uses a fixed public benchmark and an internally
held-out test partition. The focused CYP2D6 rerun passed input, artifact-hash,
and numeric checks but is exploratory because test-role metrics were examined
before protocol freeze. The broader confirmatory external matrix was held
after its outcome-blind power screen failed, and the B3DB candidate also
failed its log-loss power gate; no temporal cohort or confirmatory external
analysis is available to establish general transportability or deployment
validity. The ExpansionRx replication is limited to validation rows, was
analyzed with a post-hoc selector, and has no test-target comparison. Second,
endpoints
and external sources may share compounds or source lineage;
endpoint-level bootstrap intervals and seed summaries are descriptive
and do not represent independent assay-level sampling. Third, the
test-set reliability oracle is an after-the-fact candidate-pool
reference and is not an implementable model-selection rule. Fourth,
several practical thresholds and the combined gate were defined in code
during analysis development rather than registered prospectively. Fifth,
Chemprop used a validation-MAE checkpoint proxy on Spearman-designated
tasks because the implementation did not track Spearman for
checkpointing. Finally, the public package does not redistribute the raw
TDC rows or checkpoints, exact remote traditional-CPU package versions
were not retained in the local manifests, and exact replay requires an
authorized matching snapshot and environment.

\subsection{6. Conclusions}\label{conclusions}

Across 22 fixed TDC ADMET endpoints, validation-based family selection
and held-out reliability were frequently misaligned, particularly for
selective prediction and applicability alignment. The results support a
multi-dimensional reliability audit as a complement to conventional
validation scores. The validation-only ExpansionRx replication shows that a
multi-dimensional selector can remain stable across held-out seeds in a
source-specific development package, but it does not establish external
performance. No confirmatory external calibration/transport result
is established: the CYP2D6 diagnostic is exploratory because of prior
test-role exposure, and the broader outcome-blind power gate failed. The
overall evidence remains retrospective and scope-bounded; new, untouched,
adequately powered external cohorts are required before stronger
transportability claims.

\subsection{7. Declarations and submission
fields}\label{declarations-and-submission-fields}

The confirmed submission facts are: Huajian Jiang is the sole, first and
corresponding author; correspondence is through jian0157@e.ntu.edu.sg;
ORCID is 0009-0001-3992-2286; no funding was received; the author reports
no relevant financial or non-financial competing interests; acknowledgements
are none; and ethics approval and informed consent are not applicable because
this retrospective computational study uses public chemical and assay data
and involves no human participants, identifiable personal data, or human
materials. CRediT contributions are: Conceptualization; Data curation;
Formal analysis; Investigation; Methodology; Software; Validation;
Visualization; Writing -- original draft; and Writing -- review and editing.
The exact AI tools, versions and dates, and the final data/code release route
still require final author confirmation before submission. AI systems were not
authors and the human author remains responsible for every scientific claim,
analysis, reference, and declaration.

\subsection{References}\label{references}

\begin{thebibliography}{15}
\bibitem{huang2021tdc}
  Huang K, Fu T, Gao W, Zhao Y, Roohani Y, Leskovec J, Coley CW, Xiao C,
  Sun J, Zitnik M. Therapeutics Data Commons: Machine Learning Datasets
  and Tasks for Drug Discovery and Development. \emph{Proceedings of the
  Neural Information Processing Systems Track on Datasets and
  Benchmarks}. 2021.
\bibitem{guo2017calibration}
  Guo C, Pleiss G, Sun Y, Weinberger KQ. On Calibration of Modern Neural
  Networks. \emph{Proceedings of the 34th International Conference on
  Machine Learning}. 2017;70:1321-1330.
\bibitem{zhao2026reliability}
  Zhao D, Zhu Y, Wu Z, Wan Y, Liu X, Li S, Xu H, Hou T, Hsieh C-Y.
  Revisiting ADMET prediction reliability under real-world challenges in
  the foundation model era. \emph{Journal of Cheminformatics}.
  2026;18:95. https://doi.org/10.1186/s13321-026-01217-2.
\bibitem{kamuntavicius2025benchmarking}
  Kamuntavičius G, Paquet T, Bastas O, Šalkauskas D, Prat A, Abdel Aty
  H, Pabrinkis A, Norvaišas P, Tal R. Benchmarking ML in ADMET
  predictions: the practical impact of feature representations in
  ligand-based models. \emph{Journal of Cheminformatics}. 2025;17.
  https://doi.org/10.1186/s13321-025-01041-0.
\bibitem{zheng2026frontier}
  Zheng J, Guo C, Wang Z, Liu X, Chen H. Beyond Scaffold Splits:
  Structural-Frontier Evaluation Reveals Hidden Failures in ADMET
  Models. arXiv:2607.10729, version 4, 2026.
  https://doi.org/10.48550/arXiv.2607.10729.
\bibitem{hosni2026}
  Hosni Z, Gillet VJ, Marchese Robinson RL. Explicit Applicability Domain
  Calculations Can Help Determine When Uncertainty Estimates Are Less
  Reliable. \emph{ACS Omega}. 2026;11(3):4722-4743.
  https://doi.org/10.1021/acsomega.5c11875.
\bibitem{parrondo2026}
  Parrondo-Pizarro R, Lanini J, Rodríguez-Pérez R. Uncertainty Quantification
  in Molecular Machine Learning for Property Predictions under Data Shifts.
  \emph{Journal of Chemical Information and Modelling}. 2026;66(2):923-935.
  https://doi.org/10.1021/acs.jcim.5c02381.
\bibitem{koleiev2026}
  Koleiev I, Stratiichuk R, Shevchuk N, Melnychenko M, Nyporko A, Todoryshyn
  D, Husak V, Starosyla S, Yesylevskyy S, Nafiiev A. An End-User Audit of
  Reproducibility, Data Leakage, and Overfitting of the Top-Ranked ADMET
  Prediction Models in TDC Leaderboards. \emph{Journal of Chemical
  Information and Modelling}. 2026;66(14):8045-8058.
  https://doi.org/10.1021/acs.jcim.6c00819.
\bibitem{castellanos2026}
  Castellanos MA, MacDermott-Opeskin H, Ainsley J, Swain JA, Colby S, West D,
  et al. Assessing the State of the Art in ADMET Prediction: Lessons from the
  OpenADMET-ExpansionRx Community Blind Challenge. \emph{ChemRxiv}. 2026;
  version 3, 5 October 2026. https://doi.org/10.26434/chemrxiv.15009342/v3.
  Preprint; not peer reviewed.\bibitem{rogers2010ecfp}
  Rogers D, Hahn M. Extended-connectivity fingerprints. \emph{Journal of
  Chemical Information and Modelling}. 2010;50(5):742-754.
  https://doi.org/10.1021/ci100050t.
\bibitem{xu2019gnn}
  Xu K, Hu W, Leskovec J, Jegelka S. How Powerful are Graph Neural
  Networks? \emph{International Conference on Learning Representations}.
  2019.
\bibitem{yang2019molecular}
  Yang K, Swanson K, Jin W, Coley C, Eiden P, Gao H, Guzman-Perez A,
  Hopper T, Kelley B, Mathea M, Palmer A, Settels V, Jaakkola T, Jensen
  K, Barzilay R. Analyzing Learned Molecular Representations for
  Property Prediction. \emph{Journal of Chemical Information and
  Modelling}. 2019;59(8):3370-3388.
  https://doi.org/10.1021/acs.jcim.9b00237.
\bibitem{breiman2001rf}
  Breiman L. Random forests. \emph{Machine Learning}. 2001;45:5-32.
  https://doi.org/10.1023/A:1010933404324.
\bibitem{chen2016xgboost}
  Chen T, Guestrin C. XGBoost: A Scalable Tree Boosting System. In:
  \emph{Proceedings of the 22nd International Conference on Knowledge
  Discovery and Data Mining}. 2016:785-794.
  https://doi.org/10.1145/2939672.2939785.
\bibitem{ke2017lightgbm}
  Ke G, Meng Q, Finley T, Wang T, Chen W, Ma W, Ye Q, Liu TY. LightGBM:
  A Highly Efficient Gradient Boosting Decision Tree. \emph{Advances in
  Neural Information Processing Systems}. 2017;30.
\end{thebibliography}

\end{document}
